\section{Diskussion}\label{diskussion}

\textit{Die vollstaendige Diskussion wird ergaenzt, sobald die Grid-Search-Optimierung und
die SHAP-Analyse abgeschlossen sind. Bereits jetzt belegbar sind die folgenden methodischen
Bezugspunkte.}

\textbf{Vergleich mit Raturi et al.}
Raturi et al.\ \cite{raturi2026exploratory} fuehren die \ba{} als Metrik ein und berichten eine
\ba{} von 0{,}952 in der Mehrklassenklassifikation, jedoch mit optimierten Hyperparametern.
Die vorliegenden Ergebnisse mit Standardparametern liegen darunter (Pipeline~B: 0{,}6301),
was die Notwendigkeit der Grid-Search-Optimierung bestaetigt.
Darueber hinaus verzichten Raturi et al.\ vollstaendig auf eine Erklaerbarkeitskomponente,
was den eigenstaendigen Beitrag der vorliegenden Arbeit begruendet.

\textbf{Vergleich mit Almahaqeri et al.}
Almahaqeri et al.\ \cite{almahaqeri2026gradient} waehlen stratifiziertes Undersampling als
Strategie gegen Klassenungleichgewicht statt einer hierarchischen Klassifikationsarchitektur.
Die eigene hierarchische Architektur adressiert das Klassenungleichgewicht strukturell, ohne
Information durch Undersampling zu verlieren.

\textbf{Vergleich mit Alharby.}
Alharby~\cite{alharby2025ciciot} dokumentiert \pca{}-basierten \gb{} auf dem CICIoT2023-Datensatz.
Die eigene Ablation Study ergaenzt diese Arbeit um eine systematische Gegenuebersteflung mit
und ohne \pca{}, einschliesslich des Einflusses auf die \shap{}-Erklaerungen.

\textbf{Bewertung der SHAP-Erklaerungsqualitaet.}
\textit{Dieser Abschnitt wird nach Vorliegen der SHAP-Ergebnisse ausgefuehrt.
Die Bewertung erfolgt anhand der Metriken aus Hermosilla et al.\ \cite{hermosilla2025xai_forensic}:
Fidelitaet, Konsistenz und Jaccard-Aehnlichkeit.}

\subsection*{Offene Forschungsfragen}

Die vorliegende Arbeit belaesst drei Forschungsfragen als offen, die in zukunftigen Arbeiten
adressiert werden koennen.
Erstens die Latenzmessung und Echtzeit-Implementierung des Klassifikators fuer den praktischen
Einsatz in ressourcenbeschraenkten \iot{}-Umgebungen, die von Raturi et al.\
\cite{raturi2026exploratory} als Forschungsluecke benannt wird.
Zweitens die Robustheit gegenueber gezielten Adversarial Attacks auf das ML-System.
Drittens die Generalisierbarkeit der \shap{}-Erklaerungen auf weitere \iot{}-Datensaetze jenseits
des CICIoT2023-Benchmarks.
